% !TEX root = ../Monografia.tex
\chapter{Cronogramas do Trabalho}
\label{ap:cronogramas}

Este apêndice registra a organização temporal adotada no desenvolvimento do Trabalho de Conclusão de Curso, dividido em duas etapas complementares: \textbf{TCC 1}, voltado à fundamentação teórica, definição do problema e desenvolvimento inicial; e \textbf{TCC 2}, dedicado à implementação final, testes experimentais, análise de resultados e redação conclusiva da monografia.
O conjunto das atividades seguiu uma sequência de pesquisa, desenvolvimento e validação, com o primeiro ciclo concentrado na concepção inicial da solução e o segundo na experimentação e consolidação dos resultados.

\section{Cronograma do TCC 1}

O TCC 1 contemplou o período de \textbf{agosto a novembro de 2025} e teve como objetivo consolidar a fundamentação teórica, definir a arquitetura do sistema e apresentar um protótipo inicial funcional.
O Quadro~\ref{tab:cronograma-tcc1} apresenta a distribuição das atividades realizadas nessa etapa.

\begin{table}[htbp]
\centering
\caption{Cronograma do TCC 1 (Agosto a Novembro de 2025).}
\label{tab:cronograma-tcc1}
\small
\renewcommand{\arraystretch}{1.25}
\begin{tabular}{|p{2.7cm}|p{4.5cm}|p{6.3cm}|}
\hline
\textbf{Período} & \textbf{Atividade} & \textbf{Descrição} \\
\hline
Semanas 1--3 (04--22 ago) & Revisão Bibliográfica e Definição do Problema & Revisão sistemática sobre monitoramento de desempenho, consumo de recursos e eficiência energética em aplicativos móveis. Definição detalhada do problema de pesquisa e dos objetivos específicos. \\
\hline
Semana 4 (25--29 ago) & Consolidação da Proposta & Discussão com o orientador e entrega da versão final da \textit{Proposta de TCC}. \\
\hline
Semanas 5--6 (01--12 set) & Levantamento de Ferramentas e Tecnologias & Avaliação de bibliotecas e APIs para coleta de métricas em Flutter. Análise comparativa entre soluções existentes como \textit{DevTools}, \textit{Perfetto} e \textit{Flipper}. \\
\hline
Semanas 7--9 (15 set--03 out) & Modelagem e Arquitetura da Solução & Definição da arquitetura, fluxos de dados, diagramas de componentes e métricas a serem monitoradas. \\
\hline
Semana 10 (06 out) & Entrega do Documento de Visão & Entrega do \textit{Documento de Visão Geral do Projeto}. \\
\hline
Semanas 11--13 (07--31 out) & Implementação Inicial & Desenvolvimento dos módulos de coleta de métricas, interface básica e integração com o aplicativo de teste. Testes iniciais de funcionalidade. \\
\hline
Semanas 14--16 (03--21 nov) & Redação da Monografia e Testes Finais & Elaboração inicial da monografia, análise dos resultados parciais, documentação da metodologia e ajustes no pacote. \\
\hline
Semana 17 (24--29 nov) & Revisão e Seminário Final & Revisão geral, formatação e entrega do \textbf{Texto Inicial da Monografia}, seguida da preparação e apresentação do seminário final do TCC 1. \\
\hline
\end{tabular}
\end{table}

\section{Cronograma do TCC 2}

O TCC 2 abrangeu o período de \textbf{março a junho de 2026}, com foco na consolidação do pacote, execução dos testes experimentais, análise qualitativa dos resultados e finalização da monografia.
Essa etapa representou a consolidação prática e científica do estudo iniciado no TCC 1.

\begin{table}[htbp]
\centering
\caption{Cronograma do TCC 2 (Março a Junho de 2026).}
\label{tab:cronograma-tcc2}
\small
\renewcommand{\arraystretch}{1.25}
\begin{tabular}{|p{2.7cm}|p{4.5cm}|p{6.3cm}|}
\hline
\textbf{Período} & \textbf{Atividade} & \textbf{Descrição} \\
\hline
Semanas 1--3 (03--21 mar) & Revisão e Organização Inicial & Revisão do trabalho anterior, consolidação dos objetivos e definição dos cenários de teste e métricas de avaliação do pacote. \\
\hline
Semanas 4--6 (24 mar--11 abr) & Consolidação do Pacote & Atualização do \texttt{collector\_flutter} com módulos de análise, otimização das rotinas de coleta e implementação do painel de visualização (\textit{dashboard}) com regras heurísticas aprimoradas. \\
\hline
Semanas 7--9 (14 abr--02 mai) & Execução dos Testes Experimentais & Realização de testes controlados no pacote, observando CPU, memória, rede e sinais de desempenho. Confronto qualitativo com sinais do painel e leituras complementares por ADB. \\
\hline
Semanas 10--12 (05--23 mai) & Análise dos Resultados & Interpretação qualitativa dos dados coletados. Avaliação da coerência das métricas, do comportamento do painel e da pertinência das recomendações sugeridas pelo pacote. \\
\hline
Semanas 13--15 (26 mai--13 jun) & Consolidação, Documentação e Entrega & Organização das evidências dos experimentos, documentação da avaliação funcional e revisão textual conforme as normas da ABNT, com entrega da monografia em \textbf{07/06/2026}. \\
\hline
Semana 16 (16--18 jun) & Redação Final e Defesa & Redação final da monografia, revisão orientada, ajustes de formatação e preparação da apresentação pública do TCC 2, com defesa agendada para \textbf{18/06/2026, às 15h}. \\
\hline
\end{tabular}
\end{table}
